\documentclass[12pt]{article}
\usepackage[utf8]{inputenc}
\usepackage[T1]{fontenc}
\usepackage{hyperref}
\title{Eyak Language Resources for K-12 Curriculum Development}
\author{Compiled by Claude Sonnet 5.5 (claude-sonnet-5-5) for Andr\'as Kornai}
\date{}
\begin{document}
\maketitle

\section{Introduction}
Eyak (ISO 639-3 \texttt{eya}) is a dormant language of south-central Alaska. Its last first-language speaker, Marie Smith Jones, died in 2008. Michael Krauss documented the language, and revitalisation work by learners, notably Guillaume Leduey, continues. The resources below are those located in search results.

\section{Primary Educational Resources}
Search results describe a 2021 Pawanka Fund project aiming to make several Eyak people fluent, with a planned accessible grammar, and a 2022 e-learning programme of virtual workshops led by Leduey:
\begin{itemize}
\item \url{https://pawankafund.org/ano/2021/daxunhyuuga-eyak-language-revitalization/}
\item \url{https://pawankafund.org/?p=4540}
\end{itemize}
No ready-made school textbooks were found.

\section{Linguistic and Reference Materials}
\begin{itemize}
\item Alaska Native Languages Archive page for Eyak, with access to the Krauss dictionary and related materials: \url{https://www.alaskanativelanguages.org/eyak}
\item Eyak dictionary documentation page (UBC wiki): \url{https://wiki.ubc.ca/Documentation:RelLex/Eyak_Dictionary}
\item Eyak language overview (Wikipedia): \url{https://en.wikipedia.org/wiki/Eyak_language}
\item Eyak writing and background (Omniglot): \url{https://omniglot.com/writing/eyak.php}
\end{itemize}

\section{Religious and Community Resources}
Search summaries mention the Eyak Preservation Council as running a learner's website; its URL was not seen directly. News coverage:
\begin{itemize}
\item \url{https://www.thecordovatimes.com/2018/08/03/eyak-language-and-culture-are-growing-globally/}
\item \url{https://www.adn.com/rural-alaska/article/reviving-dead-alaska-native-language/2011/06/23}
\end{itemize}

\section{Digital and Audio Resources}
Digital copies of the dictionary and drafts are reported as open access via the Alaska Native Language Archive (see above). No separate audio collection was confirmed in the results.

\section{Curriculum Development Context}
There are no first-language speakers, so a school curriculum would rest on archival documentation and a small community of learners and teachers. No school programme was found in the results.

\end{document}